% Generated draft: requires mathematical and source review.
\documentclass[11pt]{article}
\usepackage[margin=25mm]{geometry}
\usepackage{fontspec}
\setmainfont{Latin Modern Roman}
\usepackage{amsmath,amssymb,mathtools,unicode-math}
\setmathfont{Latin Modern Math}
\usepackage{enumitem,needspace,xurl,hyperref}
\setlength{\parindent}{0pt}
\setlength{\parskip}{5pt}
\setlength{\emergencystretch}{3em}
\newcommand{\N}{\mathbb{N}}
\newcommand{\Z}{\mathbb{Z}}
\newcommand{\Q}{\mathbb{Q}}
\newcommand{\R}{\mathbb{R}}
\newcommand{\C}{\mathbb{C}}
\newcommand{\F}{\mathbb{F}}
\newcommand{\A}{\mathbb{A}}
\newcommand{\PP}{\mathbb{P}}
\newcommand{\E}{\mathbb{E}}
\newcommand{\eps}{\varepsilon}
\newcommand{\dd}{\,\mathrm{d}}
\begin{document}
\section*{UnsolvedMath ID: 700748}
\textbf{Source problem code:} catalogue entry
\subsection{Definitions and mathematical statement}
A group $G$ is \emph{residually finite} if, for every $g\in G\setminus\{1\}$, there exist a finite group $F$ and a homomorphism $\varphi:G\to F$ such that $\varphi(g)\ne1$. A subgroup $A\le G$ is abelian if $ab=ba$ for all $a,b\in A$. For $g\in G$, its centralizer is
\[
C_G(g)=\{x\in G:xg=gx\}.
\]
\textbf{(a)} Is it true that every infinite residually finite group $G$ has an infinite abelian subgroup? Explicitly, does
\[
\forall G\;\bigl[(G\text{ is infinite and residually finite})\Longrightarrow \exists A\le G\;(A\text{ is infinite and abelian})\bigr]
\]
hold? The supplied source states that this is equivalent to asking whether
\[
\forall G\;\bigl[(G\text{ is infinite and residually finite})\Longrightarrow \exists g\in G\setminus\{1\}\;|C_G(g)|=\infty\bigr].
\]
This equivalence is recorded as a source assertion.

The source attributes to Shunkov the assertion that a torsion group containing an involution with finite centralizer is virtually soluble. Here a torsion group is a group whose every element has finite order, an involution is an element of order $2$, and virtually soluble means having a soluble subgroup of finite index. A group $H$ is soluble if its derived series, defined by $H^{(0)}=H$ and $H^{(n+1)}=[H^{(n)},H^{(n)}]$, reaches the trivial subgroup after finitely many steps. The source consequently says that one may assume all elements have odd order. This is supplied background motivating a reduction; it does not restrict the universal quantifier in part (a), and its justification and the precise hypotheses of the attributed theorem require review.

\textbf{(b)} The source next asks whether every infinite ``residually $p$-group'' contains an infinite abelian subgroup. The residual class and the torsion convention are not defined in the supplied record. For precision, let $p$ be a prime and say that $G$ is \emph{residually finite $p$} if every $g\ne1$ survives under some homomorphism from $G$ to a finite group of order a power of $p$. If this is the intended meaning, part (b) asks, for each prime $p$, whether
\[
\forall G\;\bigl[(G\text{ is infinite and residually finite }p)\Longrightarrow \exists A\le G\;(A\text{ is infinite and abelian})\bigr].
\]
If instead the intended class consists of residually finite groups that are also torsion $p$-groups, the hypothesis must additionally require
\[
\forall g\in G\;\exists n\in\Z_{\ge0}\quad g^{p^n}=1.
\]
These interpretations are not identified here. The source does not explicitly restrict $p$ to odd primes; its preceding odd-order reduction is retained as background rather than silently imposed on part (b).
\subsection{Short English statement}
A. Mann asks whether every infinite residually finite group has an infinite abelian subgroup. The source states an equivalent question about an infinite centralizer of a nonidentity element, cites a Shunkov theorem to motivate considering odd-order torsion groups, and proposes a related question for ``residually $p$-groups.'' The exact meaning of that last term requires clarification. The registry labels the problem open; this is supplied status information.
\subsection{Sources}
\begin{itemize}
\item Local UnsolvedMath registry, permanent ID $700748$, exact source problem code \texttt{catalogue entry}, title \emph{Kourovka Notebook Problem 11.56}. Supplied upstream revision: \texttt{d371ee917ac9647fbb06cd6fc529deec6f3dda82}. The record was updated on 2026-10-16 and labels the problem \emph{open}; that label is not an independently verified literature assessment.
\item \emph{The Kourovka Notebook (Unsolved Problems in Group Theory)}, original problem 11.56, 11th issue (1990); proposed by A. Mann. The supplied background also identifies the 21st issue edition (2026), \texttt{arXiv:1401.0300v48}. The registry collection names E. I. Khukhro and V. D. Mazurov as editors.
\item The supplied statement attributes a theorem about torsion groups with an involution of finite centralizer to Shunkov, but supplies no separate bibliographic reference for that theorem. Its hypotheses and the claimed reduction have not been independently checked.
\end{itemize}
\end{document}
